\documentclass[12pt]{article}

\usepackage{amsmath,amssymb,amsfonts}
\usepackage{eurosym}
\usepackage[margin=1in,top=1in,bottom=1in]{geometry}

\usepackage{booktabs}
\usepackage{tabularx}
\usepackage{threeparttable}
\usepackage{array}
\usepackage{makecell}

\usepackage{graphicx}
\usepackage{caption}
\usepackage{subcaption} % preferred over subfigure
\usepackage[section]{placeins}
\usepackage{pdflscape}

%\usepackage{newtxtext}
%\usepackage{newtxmath}

\usepackage{soul,xcolor}
\usepackage[normalem]{ulem}
\usepackage{setspace}
\usepackage{sectsty}
\usepackage{comment}
\usepackage{footmisc}
\usepackage[authoryear]{natbib}
\usepackage{hyperref}
\hypersetup{
    colorlinks=true,
    linkcolor=blue,
    citecolor=blue,     
    urlcolor=cyan,
    }
    
\normalem

\onehalfspacing
\newtheorem{theorem}{Theorem}
\newtheorem{corollary}[theorem]{Corollary}
\newtheorem{proposition}{Proposition}
\newenvironment{proof}[1][Proof]{\noindent\textbf{#1.} }{\ \rule{0.5em}{0.5em}}

\newtheorem{hyp}{Hypothesis}
\newtheorem{subhyp}{Hypothesis}[hyp]
\renewcommand{\thesubhyp}{\thehyp\alph{subhyp}}

\newcommand{\red}[1]{{\color{red} #1}}
\newcommand{\blue}[1]{{\color{blue} #1}}

\newcolumntype{L}[1]{>{\raggedright\let\newline\\arraybackslash\hspace{0pt}}m{#1}}
\newcolumntype{C}[1]{>{\centering\let\newline\\arraybackslash\hspace{0pt}}m{#1}}
\newcolumntype{R}[1]{>{\raggedleft\let\newline\\arraybackslash\hspace{0pt}}m{#1}}

\captionsetup{labelfont=bf, font=small}
\newcolumntype{Y}{>{\centering\arraybackslash}X}
\newcommand{\sym}[1]{\rlap{$^{#1}$}}
\newcommand{\est}[3]{\makecell[c]{#1\sym{#3}\\ \footnotesize(#2)}}

\geometry{left=1.0in,right=1.0in,top=1.0in,bottom=1.0in}

\begin{document}

\begin{titlepage}
\title{Labor Input Responses to Ownership Changes: Evidence from U.S. Nursing Homes\thanks{abc}}
\author{Joseph White\thanks{abc} \and John Bowblis PhD\thanks{abc}}
\date{\today}
\maketitle
\begin{abstract}
\noindent The structure of my intro should be the following paragraphs. (1) Hook, introduce ownership changes and nursing homes. Introduce research question (2) Why do nursing homes matter and why are they good as a way to study this. (3) Why is staffing the main input to nursing home production. (4) Why is staffing so important to patients and quality of care. (5) Two types of ownership changes, talk about type 1. (6) Type 2, this is the type we study. (7 \& 8) Conceptual framework and the cost quality trade off. (9) Measurement, Restate Thesis Question and maybe Hypothesis, and then empirical strategy/preview findings.  \\
\vspace{0in}\\
\noindent\textbf{Keywords:} key1, key2, key3\\
\vspace{0in}\\
\noindent\textbf{JEL Codes:} key1, key2, key3\\

\bigskip
\end{abstract}
\setcounter{page}{0}
\thispagestyle{empty}
\end{titlepage}
\pagebreak \newpage




\doublespacing


\section{Introduction} \label{sec:introduction}

Ownership changes are increasingly common in healthcare markets and have become a growing concern for policymakers. In U.S. Nursing homes, Skilled nursing facilities (SNFs) can change ownership through acquisitions, sales, and corporate reorganizations, and these changes are frequently talked about because of their possible implications to the quality of care (SOURCES). Between 2016 and 2021 more than 3,200 SNFs were sold in the United States \citep{prusynski_humbert_mroz_2023}, and this number will continue to increase well into the future (SOURCE). Nursing homes serve medically vulnerable residents who depend on staff for both oversight and daily living assistance, even small changes in staffing could translate into significant differences in care. We study nurse staffing as a central input to the nursing home care production, through which ownership changes may affect care delivery. Specifically we ask, how do ownership changes affect nurse staffing levels in U.S. nursing homes, and do effects differ in chains vs non-chains and before versus during the COVID-19 pandemic?

Nursing homes provide a particularly important setting in which to study the non-price consequences of ownership changes. Shortcoming in nursing home quality has been a long standing concern, and the policy debate has repeatedly emphasized incentives and constraints that shape providers’ quality choices \citep{hackmann_2019}. Nursing home markets also exhibit severe frictions: residents and families place substantial weight on proximity, making facility decisions central to the care experienced by local populations \citep{gupta_et_al_2021}. These frictions imply that ownership changes can affect resident welfare through internal production decisions even when the facility remains open and continues to serve the same people facing the same regulations. Staffing is an especially natural outcome to study because it is both highly important for quality and one of the few outcomes that can adjust relatively quickly in response to changes in ownership.

Nursing labor is the core input in nursing home care production, and staffing is one of the most consequential operation choices facilities make.We distinguish among registered nurses (RNs), licensed practical nurses (LPNs)and certified nursing assistants (CNAs). RNs include registered nurses and directors of nursing and play central supervisory and clinical roles, including assessing residents’ conditions, developing and coordinating care plans, overseeing medication-related processes, and supervising other staff. LPNs provide routine clinical services and medication management tasks such as monitoring vital signs, while CNAs deliver the majority of hands-on daily care, assisting with bathing, dressing, toileting, feeding, and mobility \citep{lin_2014}. Because these staff types differ in training, scope of practice, and cost, ownership changes may affect not only total staffing hours but also the skill mix of care delivered. For example, an owner seeking quick cost containment may reduce total hours or substitute away from higher-cost licensed nurses, while an owner focused on compliance or clinical quality may increase RN oversight or adjust supervision structures.


A large body of evidence indicates that staffing is directly linked to nursing home quality. Early work shows that fewer RN hours and fewer aide hours are associated with more deficiencies, including deficiencies related to quality of care and quality of life \citep{harrington_et_al_2000}. Importantly, however, staffing is an endogenous choice made jointly with other inputs and under a set of constraints such as regulation and budgets that complicate causal interpretation of simple staffing and quality correlations. 

Paragraph on Evidence for why staffing is so important to quality

Our focus on staffing responses to ownership changes is also motivated by literature which shows that financial incentives and ownership structure can have an influence care in nursing homes. Private Equity (PE) ownership of nursing homes has been a rising policy concern, and PE ownership incentives can affect patient outcomes and operational decisions.
\cite{gupta_et_al_2021} find large increases in short-term mortality associated with PE owned nursing homes and show that staffing composition changes following PE acquisition. They find declines in LPNs and CNAs alongside increases in RN hours, highlighting staffing as a mechanism linking ownership incentives to resident outcomes. While private equity represents one specific ownership form, these findings prove a more general point that is central to this paper: changes in ownership can shift objectives and constraints in ways that reshape internal production decisions, including both staffing levels and skill mix.

\section{Data and Methods}
\subsection{Data Sources}
We utilize three main sources of data: Nursing Home Compare (NHC) Archive data, Healthcare Cost Report Information System (HCRIS), and Payroll Based Journal (PBJ) data. The NHC Archive contains monthly facility level information on ownership, chain affiliation, certification status, and quality measures. The second is HCRIS, also known as the Medicare Cost Reports, which contains annual cost reports submitted by every Medicare certified nursing home in the United States. The cost reports contain information about costs, revenue, financial data for each reporting facility, and dates of ownership changes. Lastly, we use the PBJ data to obtain staffing information. The PBJ data contains the census and the number of hours worked by various types of NH staff on a daily basis.

\subsection{Sample selection}

Our sample is constructed at the facility-year-month level. We create our sample using the following criteria. First, we limit our sample to the 48 contiguous U.S. states from the period January 2017 to June 2024. Secondly, the NHC Archive data is matched to the HCRIS data on a unique nursing home provider number. Because HCRIS only contains information on freestanding nursing homes, our sample excludes hospital-based nursing homes. Hospital based nursing homes differ from freestanding facilities in their structure and financing. Because they are within hospitals, staffing policies and resource allocation may be jointly determined with hospital operations and respond to hospital shocks and regulations rather than nursing home specific incentives. We therefore exclude hospital based facilities to maintain a more homogeneous set of providers and to ensure ownership transitions reflect changes in nursing home management rather than broader hospital system changes.  Hospital based nursing homes make up about 12 percent of all nursing homes in the United States. 

%Lastly, if the facilities in our sample have more than one ownership change during our sample period then we only focus on the first change \textbf{ASK DR BOWBLIS ABOUT THIS}

\subsection{Identifying Ownership Changes}

No single data source provides a comprehensive record of ownership transitions for U.S. nursing homes, so we combine information from HCRIS and the NHC Archive in a two-step identification and verification process to identify ownership changes.

The first step was to use the HCRIS data which contains an indicator for an ownership change, and corresponding date, in each cost report after an ownership change. We use this as our flag for ownership changes during the study period. However, these indicators may capture administrative or legal changes that do not reflect meaningful transfers of control. Thus, to verify these changes, our second step was to validate each potential ownership change in the HRCIS data with the NHC Archive data. The NHC Archive data reports ownership shares for each owner associated with a facility on a monthly basis. One issue with the NHC Archive data is that it reports the direct owners and indirect owners of each NH. For example, a limited liability corporation may directly own the NH, but that limited liability corporation is owned by different individuals. Therefore, we identified ownership changes to occur among direct and indirect owners using the following criteria: 

Using this information, we define a change in ownership when a majority of ownership shares turn over between month $t-1$ and month $t$. Intuitively, this captures transactions in which control of the facility changes hands, including cases where a new operator acquires most or all shares. Secondly, when ownership percentages are partially or completely missing, we approximate shares by normalizing any reported percentages and, where necessary, assigning equal weights across remaining owners. To reduce false positives from purely legal restructuring, we do not flag a change in ownership if at least 80 percent of the ownership stake remains within the same surname family (e.g., a change from “Smith Holdings LLC” to “Smith Family Partners LP”). 

Using these verified ownership changes, the analytic sample included Nursing Homes that either (i) have no ownership change in either source or (ii) have exactly one ownership change in each source with the ownership change date within six months of one another. For these Nursing Homes, we define the change in ownership month as the first month in which the NHC Archive data indicate a change. All dates are converted to calendar months, so changes occurring at the beginning or end of a month are treated identically. 

\subsection{Nursing Staff Levels}

Our main dependent variables are measures of nursing staff levels. The nursing staff types we use in our study are registered nurses (RNs), licensed practical nurses (LPNs), and certified nursing assistants (CNAs). RNs are the most highly trained nursing staff and are responsible for assessment, care planning, medication administration, and supervision of other nursing personnel. LPNs provide basic medical care under the supervision of RNs and physicians, including monitoring vital signs and administering certain treatments. CNAs provide the majority of direct daily care to residents, assisting with activities of daily living. Because these staff types differ in training, scope of practice, and compensation, changes in staffing composition may reflect different organizational responses to ownership changes. We therefore examine each staff category separately, as well as total nursing hours per patient day. 

Nursing staff levels were calculated from PBJ data, which reports the number of hours worked and the census for each type of nursing staff on a daily basis. We used this information to calculate nursing staff levels in hours per patient day (HPPD) as follows: for each facility $i$, calendar month $t$, and staff type, we first aggregate daily hours to the monthly level and then compute
\begin{equation}
  \text{HPPD}_{i,t} = \frac{H_{i,t}}{R_{i,t}},
\end{equation}
where $H_{i,t}$ is the total hours worked by the staff type during month $t$ and $R_{i,t}$ is the total resident days in that facility-month. Total HPPD is defined as the sum of RN, LPN, and CNA HPPD.

We construct two additional measures of data quality. First, we define a monthly coverage ratio
\[
  C_{i,t} = \frac{\text{DaysReported}_{i,t}}{\text{DaysInMonth}_t},
\]
where $\text{DaysReported}_{i,t}$ is the number of distinct days in month $t$ with complete PBJ records for facility $i$. This measure allows us to identify incomplete data in our sample and describe our data completeness. We use this measure for sample construction and descriptive purposes, excluding facilities and observations with low reporting coverage.

Second, we track gaps in the reporting for each facility. For each observation we compute \textit{gap}, defined as the number of calendar months since the previous PBJ observation for that facility. This captures discontinuities in reporting that may arise from administrative faults or temporary periods of not reporting PBJ system, for example during the COVID-19 period. Our baseline specification places no restriction on reporting gaps, however; we use the gap measure in robustness analyses to assess whether our results are sensitive to excluding observations following reporting gaps.

Finally, we exclude a small number of observations with implausible staffing levels: months in which RN and LPN HPPD are both zero, total HPPD is below 1.5 or above 12, or CNA HPPD exceeds 5.25. These observations are outliers that likely reflect reporting errors rather than genuine variation in staffing. These exclusion criteria align with the criteria outlined in the CMS NHC Technical Users Guide.

\subsection{Other Controls}

To account for other factors that are correlated with nursing staff levels, we included a number of control variables that are consistent with other studies (XXX). These control variables were measured at the facility-level and came from NHC and HCRIS. The following variables are used as controls in all specifications: Ownership type (government, non-profit, and for-profit), chain affiliation, total number of beds, occupancy rate, payer mix (percent of revenue coming from Medicare, and percent of revenue coming from Medicaid), and acuity. Our measure of acuity is case-mix provided by CMS in the NHC archive data. During our study period CMS altered the methodology used to construct case-mix indices. As a result, raw case-mix values are not directly comparable over time. To address this issue, we measure acuity using relative rankings rather than levels. Specifically, we construct case-mix quartiles within each state and calendar month.

\subsection{Analytic Approach}

If nursing homes changed at the same time we would use classic diff in diff, but since they don't, we dont. In our empirical setting, nursing homes change ownership at different points in time. Recent literature shows that conventional two-way fixed effects (TWFE) estimators can produce biased estimates in staggered adoption settings when treatment effects are heterogeneous over time \citep{goodman_bacon_2021}. In particular, TWFE estimates can place negative weights on certain comparisons and implicitly rely on already treated units as controls. To circumvent this, we employ a staggered adoption event study regression design. This approach allows us to estimate treatment effects before and after ownership changes while assessing pre-trends. Specifically we estimate the following event study regression:

\begin{equation}
\label{eq:es_twfe}
Y_{it}
=
\sum_{\tau = -24}^{24}
    \beta_{\tau}\,
    \mathbf{1}\!\left\{ \text{event\_time}_{it} = \tau \right\}
    \cdot \text{Treat}_{i}
+
X_{it}^{\prime}\gamma
+
\alpha_i
+
\lambda_t
+
\varepsilon_{it}.
\end{equation}

where $Y_{it}$ denotes nursing staff levels in HPPD for facility $i$ in in month $t$. The variable $\text{event\_time}_{it}$ measures the number of months relative to the month of ownership change. Thus, the coefficients $\beta_{\tau}$ capture the evolution of the nursing staff levels before and after ownership transitions. The vector $X_{it}$ includes time-varying facility characteristics, while $\alpha_i$ and $\lambda_t$ denote facility and month fixed effects, respectively. Identification comes from within-facility changes in staffing relative to the timing of ownership transitions, using facilities that have not yet experienced a change as controls.

In addition to estimating using levels, we also estimate specifications where the dependent variable is expressed in logarithms, log($Y_{it}$). Log specifications allow coefficients to be interpreted as approximate percentage changes in staffing and reduce sensitivity to differences in baseline staffing levels across facilities. For outcomes with zero staffing in a given month, log specifications are omitted.

\subsubsection{Parallel Trends}

The identifying assumption underlying our event-study design is that, absent an ownership change, staffing outcomes in treated facilities would have followed the same trends as those in untreated facilities. We assess the plausibility of this assumption by examining estimated coefficients on the pre-treatment event-time indicators and by conducting joint tests of pre-trends. Specifically, we test whether the coefficients on all pre-treatment leads of the ownership change are jointly equal to zero using Wald tests. 

Table~\ref{tab:pretrend-pvals-by-tau} reports p-values for individual pre-treatment event-time coefficients across outcomes, separately for specifications that include all observations (Panel~A) and for specifications that exclude the anticipatory period immediately preceding ownership changes (Panel~B). Under the parallel trends assumption, we would expect pre-treatment coefficients to be statistically indistinguishable from zero, implying relatively large p-values and no systematic pattern of significance across pre-treatment periods.

\begingroup
\begin{table}[!ht]
\centering
\begin{threeparttable}
\caption{Joint Wald Tests of Pre-trends (Event Study)}
\label{tab:pretrend-wald-tests}
\small
\setlength{\tabcolsep}{6pt}

\begin{tabularx}{\textwidth}{@{} l YYYY @{} }
\toprule
 & \multicolumn{4}{c}{\textbf{Outcomes}} \\
\cmidrule(lr){2-5}
 & \textbf{RN} & \textbf{LPN} & \textbf{CNA} & \textbf{Total} \\
\midrule
\multicolumn{5}{@{}l}{\textbf{Panel A: Levels (HPPD)}} \\[2pt]
With anticipation & $46.42$ (11) [0.0000] & $12.17$ (11) [0.3512] & $16.98$ (11) [0.1083] & $20.35$ (11) [0.0408] \\
Without anticipation & $7.55$ (8) [0.4788] & $10.49$ (8) [0.2322] & $3.89$ (8) [0.8666] & $6.08$ (8) [0.6387] \\
\addlinespace[6pt]
\multicolumn{5}{@{}l}{\textbf{Panel B: Logs (HPPD)}} \\[2pt]
With anticipation & $29.02$ (11) [0.0023] & $17.71$ (11) [0.0886] & $33.75$ (11) [0.0004] & $23.95$ (11) [0.0129] \\
Without anticipation & $3.33$ (8) [0.9122] & $13.48$ (8) [0.0965] & $7.50$ (8) [0.4837] & $6.33$ (8) [0.6102] \\
\bottomrule
\end{tabularx}

\begin{tablenotes}[flushleft]
\footnotesize
\item \textit{Notes:} Each cell reports the Wald $\chi^2$ statistic for the joint null that all pre-treatment event-time coefficients equal zero, followed by degrees of freedom in parentheses and the p-value in brackets.
\item Tested windows: With anticipation tests $\tau=-12$ to $\tau=-2$; Without anticipation tests $\tau=-12$ to $\tau=-5$ (since $\tau=-3,-2,-1$ are dropped and $\tau=-4$ is the reference).
\item Reference periods: With anticipation uses $\tau=-1$; Without anticipation uses $\tau=-4$.
\item Sample sizes (rows): With anticipation ($N=632,231$); Without anticipation ($N=628,107$).
\item All specifications include facility and month fixed effects and covariates: \textit{government}, \textit{non-profit}, \textit{chain}, \textit{beds}, \textit{occupancy rate}, \textit{percent Medicare}, \textit{percent Medicaid}, and state case-mix quartile indicators.
\end{tablenotes}
\end{threeparttable}
\end{table}
\endgroup


\FloatBarrier

In the ``2 Year Full Pre-Window'' specification, which includes the full set of pre-treatment months ($\tau=-24$ to $\tau=-2$) and uses $\tau=-1$ as the reference period, the joint Wald tests reject the null of no pre-trends for several outcomes, particularly for RN and total staffing in both levels and logs. This indicates that staffing shows change effects in the periods leading up to the recorded ownership transition when the months immediately preceding the change are included in the pre-window, consistent with anticipatory adjustments around the transition date.

We therefore also report ``donut'' specifications that exclude the three months immediately preceding ownership changes ($\tau\in\{-3,-2,-1\}$) and use $\tau=-4$ as the reference period. Relative to the full pre-window tests, the donut specifications yield substantially weaker evidence of pre-trend violations for several outcomes, with joint tests that are no longer statistically distinguishable from zero in the shorter one-year donut window. Overall, these results suggest that the months immediately before the recorded change month are likely contaminated by anticipation, and that excluding $\tau\in\{-3,-2,-1\}$ improves the plausibility of the parallel trends assumption.

Accordingly, our preferred specifications estimate equation~\eqref{eq:es_twfe} excluding observations in $\tau\in\{-3,-2,-1\}$ and interpret dynamic treatment effects relative to $\tau=-4$. Standard errors are clustered at the facility level to account for serial correlation within nursing homes.

While our main specification relies on the event-study design, we also present TWFE estimates to summarize the average post ownership change effect. We estimate the following regression:

\begin{equation}
  Y_{it}
  =
  \beta_{\text{post}}\,\textit{Post}_{it}
  +
  X_{it}^{\prime}\boldsymbol{\gamma}
  +
  \alpha_i
  +
  \lambda_t
  +
  \varepsilon_{it},
  \label{eq:twfe-post}
\end{equation}

where $Y_{it}$ denotes nurse staffing outcomes for staff types RN, LPN, CNA, or total HPPD for facility $i$ in month $t$. The indicator $\textit{Post}_{it}$ equals one for all months following an ownership change and zero otherwise. Facility fixed effects $\alpha_i$ absorb time invariant differences across nursing homes, while month fixed effects $\lambda_t$ control for common shocks to staffing over time.

The coefficient $\beta_{\text{post}}$ captures the average change in staffing following an ownership transition, relative to the pre-change period. All specifications include the same time varying covariates as in the event-study models, and standard errors are clustered at the facility level.

\subsubsection{Subgroup Analysis}

We further examine heterogeneity in the effects of ownership change with two further analyses: the COVID-19 pandemic and baseline chain affiliation.

First, we estimate separate models for the pre-pandemic period (January 2017 through December 2019) and the pandemic period (April 2020 through June 2024). Labor market conditions for nursing staff changed substantially during the COVID-19 pandemic, including heightened turnover, changes in wages, and staffing shortages. Estimating models separately by period allows us to assess whether ownership changes had differential effects on staffing before and during this structural shift in the nursing labor market.

Second, we examine heterogeneity by baseline chain status. We classify facilities based on whether they were affiliated with a chain at the beginning of the study period (January 2017), prior to any ownership transitions in our sample. Chain affiliated facilities differ systematically from independent facilities in ways that are directly relevant for staffing decisions. Prior work shows that chains operate internal labor markets, share managerial expertise across facilities, and are better positioned to reallocate staff in response to shocks, potentially lessening the effects of ownership changes on staffing outcomes. In contrast, independent facilities may face tighter local labor constraints and fewer organizational buffers, making staffing levels more sensitive to changes in ownership and management.

By looking at baseline chain status, we compare facilities with similar organizational structures prior to treatment, rather than confounding the effects of ownership change with changes in chain affiliation occurring at the same time. Estimating models separately for chain and non-chain facilities therefore allows us to assess whether ownership transitions have different effects on staffing across organizational structures, and whether chain affiliation dampens or magnifies staffing adjustments following ownership changes.

In both subgroup analyses, we estimate the same specifications as in equations~\eqref{eq:es_twfe} and~\eqref{eq:twfe-post}, maintaining consistent controls, fixed effects, and clustering. These analyses shed light on mechanisms and contextual factors that shape the staffing consequences of ownership changes.

\section{Results}

\subsection{Summary Statistics}

Table~\ref{tab:sumstats} reports summary statistics for the main analysis sample. The final sample consists of approximately 7,806 unique skilled nursing facilities and 632,231 facility-month observations spanning January 2017 through June 2024. Of these facilities,1589 experience an ownership change during the sample period.

\begin{table}[!ht]
\centering
\begin{threeparttable}
\caption{Summary Statistics}
\label{tab:sumstats}
\small
\setlength{\tabcolsep}{8pt}

\begin{tabularx}{\textwidth}{@{} l r r @{} }
\toprule
\textbf{Variable} & \textbf{Mean} & \textbf{SD} \\
\midrule
\multicolumn{3}{@{}l}{\textbf{Panel A: Outcome variables}} \\[2pt]
RN HPRD & 0.427 & 0.309 \\
LPN HPRD & 0.796 & 0.325 \\
CNA HPRD & 2.100 & 0.548 \\
Total HPRD & 3.324 & 0.793 \\
\addlinespace[0.8em]
\multicolumn{3}{@{}l}{\textbf{Panel B: Control variables}} \\[2pt]
Government (dummy) & 0.101 & 0.301 \\
Non-profit (dummy) & 0.171 & 0.377 \\
Chain affiliation (dummy) & 0.546 & 0.498 \\
Beds & 108.4 & 57.9 \\
Occupancy rate (\%) & 76.4 & 15.5 \\
\% Medicare & 13.0 & 11.1 \\
\% Medicaid & 60.5 & 20.2 \\
Acuity quartile 2 (state-month) & 0.231 & 0.422 \\
Acuity quartile 3 (state-month) & 0.230 & 0.421 \\
Acuity quartile 4 (state-month) & 0.233 & 0.423 \\
\bottomrule
\end{tabularx}

\begin{tablenotes}[flushleft]
\footnotesize
\item \textit{Notes:} The unit of observation is facility--month. HPRD denotes hours per resident day. RN, LPN, and CNA denote registered nurses, licensed practical nurses, and certified nursing assistants, respectively. Total HPRD is the sum of RN, LPN, and CNA HPRD. Occupancy rate is the ratio of residents to certified beds (percent). \% Medicare and \% Medicaid are the shares of residents covered by Medicare and Medicaid, respectively. Government and Non-profit are ownership-type indicator variables (for-profit omitted category). Chain affiliation is an indicator for membership in a multi-facility chain. Acuity quartiles are state--month quartiles of resident acuity (case-mix) with quartile 1 omitted. Rows $=$ 632,231; Facilities $=$ 7,806; Treated facilities $=$ 1,589; Period $=$ 2017/01–2024/06; Average months per facility $=$ 81.0.
\end{tablenotes}

\end{threeparttable}
\end{table}

\FloatBarrier

\subsection{Main Results}

Given our main analysis suggests that parallel trends assumption likely holds, Figure~\ref{fig:event_study_main_landscape} presents event study estimates of the effect of ownership changes on nursing staffing outcomes (i.e., Equation~\ref{eq:es_twfe}). Following an ownership transition, staffing declines persist for many periods, with the largest and most precise effects occurring for registered nurses (RN), certified nursing assistants (CNA), and total nursing hours per patient day.

RN staffing declines immediately following ownership changes and remains far below pre-transition levels throughout the post-treatment period. CNA staffing exhibits a similar but slightly smaller decline, while total nursing hours mirror the combined reductions across staff types. In contrast, LPN staffing shows limited and statistically insignificant changes over time. These dynamic patterns indicate that ownership transitions lead to sustained reductions in staffing rather than short run adjustments.


\newgeometry{left=1.2cm,right=1.2cm,top=1.5cm,bottom=1.5cm}

\begin{landscape}
\begin{figure}[p]
\centering

% Row 1
\begin{subfigure}[t]{0.49\linewidth}
  \centering
  \includegraphics[width=\linewidth]{C:/Repositories/white-bowblis-nhmc/outputs/plots/twfe_es_rn_without_anticipation_drop_3_1_.pdf}
  \caption{Registered Nurses (RN)}
\end{subfigure}\hfill
\begin{subfigure}[t]{0.49\linewidth}
  \centering
  \includegraphics[width=\linewidth]{C:/Repositories/white-bowblis-nhmc/outputs/plots/twfe_es_lpn_without_anticipation_drop_3_1_.pdf}
  \caption{Licensed Practical Nurses (LPN)}
\end{subfigure}

\vspace{1em}

% Row 2
\begin{subfigure}[t]{0.49\linewidth}
  \centering
  \includegraphics[width=\linewidth]{C:/Repositories/white-bowblis-nhmc/outputs/plots/twfe_es_cna_without_anticipation_drop_3_1_.pdf}
  \caption{Certified Nursing Assistants (CNA)}
\end{subfigure}\hfill
\begin{subfigure}[t]{0.49\linewidth}
  \centering
  \includegraphics[width=\linewidth]{C:/Repositories/white-bowblis-nhmc/outputs/plots/twfe_es_total_without_anticipation_drop_3_1_.pdf}
  \caption{Total Nursing Hours}
\end{subfigure}

\caption{
Dynamic Effects of Ownership Change on Nursing Staffing.
Event study estimates of the effect of a change in ownership on nursing staff levels in hours per patient day (HPPD).
Points show coefficient estimates relative to the month immediately preceding the ownership change;
vertical bars denote 95\% confidence intervals.
The three months prior to the ownership change ($\tau=-3,-2,-1$) are excluded.
All specifications include control variables and facility and month fixed effects.
}
\label{fig:event_study_main_landscape}

\end{figure}
\end{landscape}
\FloatBarrier
\restoregeometry

\normalsize
\doublespacing

To summarize these dynamic effects with a single average treatment estimate, Table~\ref{tab:twfe-post-full} reports two-way fixed effects estimates of the ownership change effect (i.e., Equation~\ref{eq:twfe-post}). Panel A reports results in levels, while Panel B reports logged results. Consistent with the event study results, ownership changes lead to meaningful reductions in staffing levels. RN staffing declines by approximately 0.03 hours per patient day (roughly 9 percent), CNA staffing falls by about 0.05 hours (approximately 3 percent), and total staffing decreases by 0.08 hours per patient day (roughly 3 percent). LPN staffing shows no statistically significant change. Estimates are nearly identical whether the observations for the three month period before the ownership change occurred is included or excluded from the model.

\begingroup
\begin{table}[!ht]
\centering
\begin{threeparttable}
\caption{Two-Way Fixed Effects Estimates of \textit{post} on Staffing Outcomes (Baseline, Without anticipation)}
\label{tab:twfe-post-full}
\small
\setlength{\tabcolsep}{6pt}

\begin{tabularx}{\textwidth}{@{} l YYYY @{} }
\toprule
 & \multicolumn{4}{c}{\textbf{Outcomes}} \\
\cmidrule(lr){2-5}
 & \textbf{RN} & \textbf{LPN} & \textbf{CNA} & \textbf{Total} \\
\midrule
HPRD & \est{$-0.032$}{$0.005$}{***}  &  \est{$\phantom{-}0.000$}{$0.006$}{}  &  \est{$-0.054$}{$0.009$}{***}  &  \est{$-0.086$}{$0.013$}{***} \\
\addlinespace[3pt]
Log(HPRD) & \est{$-0.092$}{$0.016$}{***}  &  \est{$\phantom{-}0.003$}{$0.009$}{}  &  \est{$-0.028$}{$0.005$}{***}  &  \est{$-0.026$}{$0.004$}{***} \\
\bottomrule
\end{tabularx}

\begin{tablenotes}[flushleft]
\footnotesize
\item \textit{Notes:} Each cell reports the coefficient on \textit{post} with two-way clustered standard errors (by facility and month) in parentheses. The table reports staffing levels (HPRD) and log staffing levels (Log(HPRD)). Sample: Without anticipation ($N_{\mathrm{HPRD}}=628,107;\ N_{\mathrm{Log}}=628,107$).
\item All specifications include facility and month fixed effects and covariates: \textit{government}, \textit{non-profit}, \textit{chain}, \textit{beds}, \textit{occupancy rate}, \textit{percent Medicare}, \textit{percent Medicaid}, and state case-mix quartile indicators.
\item Statistical significance: $^{***}p<0.01$, $^{**}p<0.05$, $^{*}p<0.10$.
\end{tablenotes}
\end{threeparttable}
\end{table}
\endgroup


\FloatBarrier

\subsection{Effects Before and During Covid-19 Pandemic}

Table~\ref{tab:twfe-prepost} examines whether these effects differ before and after the COVID-19 pandemic. The magnitude of the decline is roughly halved in the post pandemic period for RN, CNA, and total HPPD, suggesting that pandemic induced shifts in labor markets and staffing substantially reduced the impact of ownership transitions.

\begingroup
\begin{table}[!ht]
\centering
\begin{threeparttable}
\caption{TWFE Estimates of \textit{post}: Pre- vs Post-pandemic Periods (Without anticipation)}
\label{tab:twfe-prepost}
\small
\setlength{\tabcolsep}{6pt}

\begin{tabularx}{\textwidth}{@{} l YYYY @{} }
\toprule
 & \multicolumn{4}{c}{\textbf{Outcomes}} \\
\cmidrule(lr){2-5}
 & \textbf{RN} & \textbf{LPN} & \textbf{CNA} & \textbf{Total} \\
\midrule
\multicolumn{5}{@{}l}{\textbf{Panel A: Staffing Levels in HPRD}} \\[2pt]
Pre-Pandemic Period (2017/01 - 2019/12) & \est{$-0.030$}{$0.007$}{***}  &  \est{$-0.023$}{$0.009$}{**}  &  \est{$-0.065$}{$0.015$}{***}  &  \est{$-0.118$}{$0.020$}{***} \\
Pandemic Period (2020/04 - 2024/06) & \est{$-0.020$}{$0.007$}{**}  &  \est{$\phantom{-}0.000$}{$0.008$}{}  &  \est{$-0.045$}{$0.016$}{***}  &  \est{$-0.064$}{$0.021$}{***} \\

\addlinespace[3pt]
\multicolumn{5}{@{}l}{\textbf{Panel B: Log Staffing Levels in HPRD}} \\[2pt]
Pre-Pandemic Period (2017/01 - 2019/12) & \est{$-0.091$}{$0.026$}{***}  &  \est{$-0.017$}{$0.013$}{}  &  \est{$-0.034$}{$0.007$}{***}  &  \est{$-0.036$}{$0.006$}{***} \\
Pandemic Period (2020/04 - 2024/06) & \est{$-0.075$}{$0.022$}{***}  &  \est{$\phantom{-}0.001$}{$0.011$}{}  &  \est{$-0.022$}{$0.008$}{**}  &  \est{$-0.021$}{$0.006$}{***} \\
\bottomrule
\end{tabularx}

\begin{tablenotes}[flushleft]
\footnotesize
\item \textit{Notes:} Each cell reports the coefficient on \textit{post} with two-way clustered standard errors (by facility and month) in parentheses. Panel~A reports levels (HPRD); Panel~B reports logs (HPRD). Sample sizes: Row~1 ($N_{\mathrm{levels}}=222,399;\ N_{\mathrm{logs}}=222,399$), Row~2 ($N_{\mathrm{levels}}=312,471;\ N_{\mathrm{logs}}=312,471$).
\item All specifications include facility and month fixed effects and covariates: \textit{government}, \textit{non-profit}, \textit{chain}, \textit{beds}, \textit{occupancy rate}, \textit{percent Medicare}, \textit{percent Medicaid}, and state case-mix quartile indicators.
\item Statistical significance: $^{***}p<0.01$, $^{**}p<0.05$, $^{*}p<0.10$.
\item Pre-pandemic 2017/01--2019/12; Pandemic 2020/04--2024/06.
\end{tablenotes}
\end{threeparttable}
\end{table}
\endgroup


\FloatBarrier

\subsection{Effects by Chain Status}

Table~\ref{tab:twfe-chain-nonchain} compares effects for facilities that began the study period as part of a chain and those that did not. Staffing reductions are noticeably larger among non chain facilities. CNA HPPD falls by 0.034 hours in chain facilities and by 0.079 hours in non chains, and total HPPD declines by 0.062 versus 0.113 hours, respectively. On a log scale, these differences translate into substantially larger percentage declines for non-chain facilities. LPN staffing again shows minimal movement.

\begingroup
\begin{table}[!ht]
\centering
\begin{threeparttable}
\caption{TWFE Estimates of \textit{post}: Chain vs Non-chain Facilities (Jan 2017 Baseline, Without anticipation)}
\label{tab:twfe-chain-nonchain}
\small
\setlength{\tabcolsep}{6pt}

\begin{tabularx}{\textwidth}{@{} l YYYY @{} }
\toprule
 & \multicolumn{4}{c}{\textbf{Outcomes}} \\
\cmidrule(lr){2-5}
 & \textbf{RN} & \textbf{LPN} & \textbf{CNA} & \textbf{Total} \\
\midrule
\multicolumn{5}{@{}l}{\textbf{Panel A: Staffing Levels in HPRD}} \\[2pt]
Chain January 2017 & \est{$-0.031$}{$0.007$}{***}  &  \est{$\phantom{-}0.002$}{$0.008$}{}  &  \est{$-0.035$}{$0.011$}{***}  &  \est{$-0.064$}{$0.015$}{***} \\
Non-chain January 2017 & \est{$-0.038$}{$0.009$}{***}  &  \est{$\phantom{-}0.003$}{$0.011$}{}  &  \est{$-0.079$}{$0.019$}{***}  &  \est{$-0.113$}{$0.026$}{***} \\

\addlinespace[3pt]
\multicolumn{5}{@{}l}{\textbf{Panel B: Log Staffing Levels in HPRD}} \\[2pt]
Chain January 2017 & \est{$-0.103$}{$0.021$}{***}  &  \est{$\phantom{-}0.005$}{$0.013$}{}  &  \est{$-0.019$}{$0.006$}{***}  &  \est{$-0.021$}{$0.005$}{***} \\
Non-chain January 2017 & \est{$-0.090$}{$0.030$}{***}  &  \est{$\phantom{-}0.015$}{$0.017$}{}  &  \est{$-0.038$}{$0.010$}{***}  &  \est{$-0.032$}{$0.008$}{***} \\
\bottomrule
\end{tabularx}

\begin{tablenotes}[flushleft]
\footnotesize
\item \textit{Notes:} Each cell reports the coefficient on \textit{post} with two-way clustered standard errors (by facility and month) in parentheses. Panel~A reports levels (HPRD); Panel~B reports logs (HPRD). Sample sizes: Row~1 ($N_{\mathrm{levels}}=318,462;\ N_{\mathrm{logs}}=318,462$), Row~2 ($N_{\mathrm{levels}}=239,781;\ N_{\mathrm{logs}}=239,781$).
\item All specifications include facility and month fixed effects and covariates: \textit{government}, \textit{non-profit}, \textit{chain}, \textit{beds}, \textit{occupancy rate}, \textit{percent Medicare}, \textit{percent Medicaid}, and state case-mix quartile indicators.
\item Statistical significance: $^{***}p<0.01$, $^{**}p<0.05$, $^{*}p<0.10$.
\item Baseline chain classification determined by facility status in January 2017.
\end{tablenotes}
\end{threeparttable}
\end{table}
\endgroup


\FloatBarrier

Overall, ownership changes lead to sizable and persistent reductions in nursing staffing, with effects concentrated among RN, CNA, and total hours per patient day. These reductions are larger prior to the COVID-19 pandemic and among facilities that were not part of a chain at baseline. Together, the event study and TWFE results indicate that ownership transitions systematically worsen staffing conditions in nursing homes rather than inducing temporary adjustments.

\subsection{Robustness}

This section presents a series of robustness checks designed to assess the validity of our identifying assumptions and the sensitivity of our results to alternative model choices. Across all exercises, the estimated effects of ownership change on staffing remain stable in sign, magnitude, and statistical significance. Most of these will have tables/plots that will go in the appendix.

\subsubsection{Assumptions}

In this section I want to present results for event study and TWFE models that rely on different assumptions than the one that I used in my main specification.


\subsubsection{Event Window and Anticipation}

Staffing decisions may adjust in anticipation of formal ownership transitions, potentially violating the parallel trends assumption if these periods are included. To address this concern, we estimate our event study and TWFE models using alternative event windows that exclude months immediately preceding the ownership change and show that the estimated post-treatment effects remain largely unchanged. Row 6 of Table~\ref{tab:twfe-robustness-summary} reports estimates of equation~\ref{eq:twfe-post} with a wider anticipatory period excluded from the sample. The results from row 6 show no differences in values or statistical significance. Similarly, row 7 reports estimates under a smaller anticipatory window that is excluded from the data, and we see no differences in the results.

\subsubsection{Gap Size}

Because staffing data are occasionally missing or reported with gaps, we examine whether our results are sensitive to the length of time between consecutive observations. Restricting the sample to facilities with smaller reporting gap shows that estimates are similar to those in the full sample, indicating that data gaps do not drive our findings. Row 2 of Table~\ref{tab:twfe-robustness-summary} reports our results for equation~\ref{eq:twfe-post} with the restriction that facilities with a reporting gap of larger than 6 are excluded from our sample. Compared to the baseline (row 1) this result is identical across staff types and statistical significance. The exact same result is true for row 3 which reports estimates for equation~\ref{eq:twfe-post} where facilities with a reporting gap of greater than 3 are removed. Row 4 and 5 report estimates where gap greater than 1 or equal to 0 respectively, are removed. Similarly, we see no differences in these results with respect to the baseline.


\section{Discussions} \label{sec:discussion}

Our findings indicate that ownership transitions in nursing homes are followed by reductions in nurse staffing, suggesting that labor inputs constitute a primary method of adjustment following changes in ownership. In a sector characterized by extensive requirements and regulatory oversight, staffing decisions represent one of the most flexible ways through which new owners can alter operating costs and organizational practices. The observed declines in RN, CNA, and total nursing HPPD point to a reallocation of resources away from labor following ownership changes. These patterns show the importance of examining internal production decisions when evaluating the consequences of ownership transitions in regulated healthcare markets.

The timing of these staffing adjustments also gives insight into the underlying effects. Event-study estimates in Figure~\ref{fig:event_study_main_landscape} show that staffing levels begin to decline shortly before the true filing transfer date and remain persistently lower in the following periods. This pattern is consistent with ownership changes being anticipated events rather than sudden shocks. During this period before the change, incumbent or incoming owners may begin restructuring operations or implementing cost-containment strategies in advance of the actual transfer date. The persistence of the reductions after ownership changes suggests that adjustments are not just transitory, but rather longer-term changes in organizational priorities and operating practices.

Labor being an adjustment margin in this context is economically intuitive because nursing labor represents one of the largest operating expenses for nursing homes and is subject to fewer regulation than prices or resident admissions. While minimum staffing requirements exist, facilities retain considerable discretion over staffing levels above the threshold, as well as over workforce composition across skill categories. In contrast, rates are largely fixed through Medicare and Medicaid, and changes in bed capacity or service lines are costly and often constrained by regulations and local market conditions. As a result, ownership changes may translate into staffing changes even when other inputs to production remain unchanged.

The heterogeneity of staffing responses across nurse types also shows how ownership changes affect labor inputs. We find larger and more consistent reductions in RN staffing relative to other nurse categories, with declines also evident among CNAs and in total staffing. This suggests that staffing changes are not uniform across skill levels and may reflect differences in wage costs. RN's require higher wages and play supervisory and clinical roles that may be more easily reallocated following ownership changes, especially in facilities that become integrated into larger organizational structures. In contrast, CNA staffing, while it still declines, may be constrained by direct care needs and minimum staffing standards. The differential response across staff types highlights how ownership transitions can alter not only staffing levels but also labor composition, with potential implications for care delivery.

The magnitude of staffing responses also differs between the pre-pandemic and pandemic periods. While ownership changes are associated with statistically significant reductions in staffing in both periods, these effects are consistently smaller during the pandemic. As shown in Table~\ref{tab:twfe-prepost}, declines in RN, CNA, and total nursing hours per patient day are larger in the pre-pandemic period than during the COVID-19 period, both in levels and in logs. One interpretation is that the pandemic fundamentally altered labor market conditions in long-term care, limiting nursing homes ability to further adjust staffing in response to ownership changes. Widespread staffing shortages of nursing labor during the pandemic may have constrained owners’ capacity to implement additional reductions. At the same time, temporary regulation and heightened scrutiny of nursing home staffing during the pandemic may have dampened the extent to which new owners could reduce labor supply.

Staffing responses also vary across chain-affiliated and non-chain facilities, indicating that organizational structure is important in shaping the labor effects of ownership changes. We find that staffing reductions following ownership changes are substantially larger, almost double in CNA and total staffing, among non-chain facilities than among chain-affiliated nursing homes. One potential explanation is that chain-affiliated facilities may benefit from greater centralization and larger labor pools. In contrast, non-chain facilities may face more binding financial constraints, making labor cost reductions a more immediate and pronounced response to changes in ownership. Additionally, independent facilities may have more limited ability to reallocate staff across locations, leading to sharper staffing reduction. These patterns suggest that ownership transitions may be more disruptive for non-chain facilities, with labor inputs bearing a much larger share of adjustment.

We have seen decreases in staffing levels in almost all  the staffing categories that we study, but an important question to ask is whether these declines matter? Do these declines in staffing actually translate to real negative affects on patients? In the baseline model, we estimate declines of $-0.032$ RN HPPD, $-0.054$ CNA HPPD, and $-0.086$ total HPPD following ownership changes (with no statistically significant change in LPN staffing). Converting these magnitudes, $-0.032$ RN HPPD corresponds to about $0.032\times 60 \approx 1.9$ fewer RN minutes per resident-day, $-0.054$ CNA HPPD corresponds to about $3.2$ fewer CNA minutes per resident-day, and $-0.086$ total HPPD corresponds to about $5.2$ fewer total nursing minutes per resident-day. While a reduction of a few minutes per resident-day may appear tiny, it aggregates to meaningful staff time at the facility level: for every 100 resident-days, a $-0.086$ total HPPD effect implies about $0.086\times 100 = 8.6$ fewer total nursing hours per day. Put differently, this is close to removing one full-time shift of nursing labor each day, sustained over time. In a setting where staffing is one of the main inputs to quality, sustained declines of this size are likely to have real implications for patient care. Figure out percentages then make comments about what this actually means. 

Our findings contribute to the literature in several ways. First, they complement the industrial organization literature on ownership changes by providing evidence on non-price outcomes in a setting where prices are mostly fixed. While prior work has emphasized price and market-power effects of mergers and acquisitions, our results highlight how ownership transitions can affect internal production decisions even in the absence of price changes. Second, this study adds to the literature on labor market outcomes following ownership changes. Previous evidence on employment and wage effects has been mixed and largely drawn from manufacturing or administrative settings, so by focusing on healthcare labor in nursing homes, we show that ownership changes can have substantial labor effects in a sector where staffing is directly tied to output quality. Third, our results contribute to the healthcare consolidation literature by identifying staffing as an important way through which ownership changes may influence care delivery, complementing existing studies that focus on prices or spending.

Paragraph on Limitations

Closing

\section{Conclusion} \label{sec:conclusion}



\singlespacing
\setlength\bibsep{0pt}
\nocite{*}
\bibliographystyle{apalike}
\bibliography{refs}

\clearpage

\onehalfspacing

\section*{Tables} \label{sec:tab}
\addcontentsline{toc}{section}{Tables}

\begingroup
\begin{table}[!ht]
\centering
\begin{threeparttable}
\caption{TWFE estimates of \textit{post} on staffing levels (HPRD): sample-restriction and case-mix robustness.}
\label{tab:twfe-robustness-summary}
\small
\setlength{\tabcolsep}{6pt}

\begin{tabularx}{\textwidth}{@{} l c YYYY @{} }
\toprule
 &  & \multicolumn{4}{c}{\textbf{Outcomes}} \\
\cmidrule(lr){3-6}
 & \textbf{N} & \textbf{RN} & \textbf{LPN} & \textbf{CNA} & \textbf{Total} \\
\midrule
(1) Baseline (no anticipation) & 550,334 & $-0.0422^{***}$ $(0.0058)$ & $\phantom{-}0.0021$ $(0.0063)$ & $-0.0641^{***}$ $(0.0099)$ & $-0.1041^{***}$ $(0.0140)$ \\
\addlinespace[0.35em]
(2) Sample excludes \textit{gap} $> 6$ & 550,175 & $-0.0422^{***}$ $(0.0058)$ & $\phantom{-}0.0018$ $(0.0063)$ & $-0.0641^{***}$ $(0.0099)$ & $-0.1045^{***}$ $(0.0140)$ \\
\addlinespace[0.35em]
(3) Sample excludes \textit{gap} $> 3$ & 549,918 & $-0.0422^{***}$ $(0.0058)$ & $\phantom{-}0.0018$ $(0.0063)$ & $-0.0641^{***}$ $(0.0099)$ & $-0.1044^{***}$ $(0.0139)$ \\
\addlinespace[0.35em]
(4) Sample excludes \textit{gap} $> 1$ & 546,319 & $-0.0424^{***}$ $(0.0058)$ & $\phantom{-}0.0022$ $(0.0063)$ & $-0.0641^{***}$ $(0.0099)$ & $-0.1042^{***}$ $(0.0139)$ \\
\addlinespace[0.35em]
(5) Sample excludes \textit{gap} $> 0$ & 546,319 & $-0.0424^{***}$ $(0.0058)$ & $\phantom{-}0.0022$ $(0.0063)$ & $-0.0641^{***}$ $(0.0099)$ & $-0.1042^{***}$ $(0.0139)$ \\
\addlinespace[0.35em]
(6) Drop $t \in \{-4,-3,-2,-1\}$ & 549,052 & $-0.0418^{***}$ $(0.0059)$ & $\phantom{-}0.0019$ $(0.0064)$ & $-0.0640^{***}$ $(0.0100)$ & $-0.1038^{***}$ $(0.0141)$ \\
\addlinespace[0.35em]
(7) Drop $t \in \{-2,-1\}$ only & 551,637 & $-0.0423^{***}$ $(0.0057)$ & $\phantom{-}0.0022$ $(0.0063)$ & $-0.0637^{***}$ $(0.0098)$ & $-0.1038^{***}$ $(0.0138)$ \\
\addlinespace[0.35em]
\addlinespace[3pt]
(8) + Case-mix: State Quartile & 550,334 & $-0.0444^{***}$ $(0.0058)$ & $-0.0002$ $(0.0064)$ & $-0.0683^{***}$ $(0.0099)$ & $-0.1128^{***}$ $(0.0139)$ \\
\addlinespace[0.35em]
(9) + Case-mix: State Decile & 550,334 & $-0.0448^{***}$ $(0.0058)$ & $-0.0004$ $(0.0064)$ & $-0.0689^{***}$ $(0.0099)$ & $-0.1142^{***}$ $(0.0139)$ \\
\addlinespace[0.35em]
(10) + Case-mix: National Quartile & 550,334 & $-0.0447^{***}$ $(0.0058)$ & $-0.0002$ $(0.0064)$ & $-0.0690^{***}$ $(0.0099)$ & $-0.1139^{***}$ $(0.0139)$ \\
\addlinespace[0.35em]
(11) + Case-mix: National Decile & 550,334 & $-0.0455^{***}$ $(0.0058)$ & $-0.0005$ $(0.0064)$ & $-0.0697^{***}$ $(0.0099)$ & $-0.1158^{***}$ $(0.0139)$ \\
\bottomrule
\end{tabularx}

\begin{tablenotes}[flushleft]
\footnotesize
\item \textit{Notes:} Each cell reports the \textit{post} coefficient on staffing levels (HPRD), with two-way clustered standard errors (facility and calendar month) in parentheses.
\item Rows (1)--(7) hold the control set fixed at Spec A (post + number of certified beds, plus facility and calendar-month fixed effects) and vary the sample restriction or anticipation-window definition, as described in each row label. Row (1) is the project's standard sample: the anticipation window ($\tau = -3, -2, -1$) excluded, no additional gap restriction.
\item Rows (8)--(11) hold the sample fixed at row (1) and add one case-mix control set on top of Spec A per row -- Spec A itself includes no case-mix control, so these rows are not directly comparable to rows (1)--(7). ``State Quartile'' is the default case-mix control used elsewhere in the paper. Quartile bins use dummies for bins 2--4; decile bins use dummies for bins 2--10 (bin 1 omitted as reference in both cases).
\item Statistical significance: $^{***}p<0.01$, $^{**}p<0.05$, $^{*}p<0.10$.
\end{tablenotes}

\end{threeparttable}
\end{table}
\endgroup


\FloatBarrier

\input{C:/Repositories/white-bowblis-nhmc/outputs/tables/pretrend_pvals_by_tau_fragment.tex}
\FloatBarrier

\clearpage

\section*{Figures} \label{sec:fig}
\addcontentsline{toc}{section}{Figures}

%\begin{figure}[hp]
%  \centering
%  \includegraphics[width=.6\textwidth]{../fig/placeholder.pdf}
%  \caption{Placeholder}
%  \label{fig:placeholder}
%\end{figure}




\clearpage

\section*{Appendix A. Placeholder} \label{sec:appendixa}
\addcontentsline{toc}{section}{Appendix A}

\end{document}